% 3
\begin{frame}[t]{LLM 与虚拟细胞的四种连接}
\vspace{.1cm}
\centering
\begin{tikzpicture}[conn/.style={box,text width=3.5cm,align=left,inner sep=7pt},
  cl/.style={font=\scriptsize,text=Muted,fill=white,inner sep=1.5pt}]
 \node[hub,text width=2.5cm,minimum height=1.6cm] (hub) at (0,0) {\ico{\Large\faProjectDiagram}\\[.1cm]\textbf{细胞预测模型}\\\small 虚拟细胞的核心};
 \node[conn,fill=PaleTeal,draw=DarkBlue,line width=.9pt] (p) at (-5,1.15) {\ico{\faLayerGroup}\enspace\textbf{借鉴预训练范式}\\[.06cm]\scriptsize Geneformer · scGPT · scFoundation\\\scriptsize\color{Muted}基因及表达量 $\rightarrow$ 细胞表征};
 \node[conn] (m) at (-5,-1.15) {\ico{\faDna}\enspace\textbf{分子与功能先验}\\[.06cm]\scriptsize UCE · Evo 2 · GenePT\\\scriptsize\color{Muted}序列或基因描述 $\rightarrow$ 表示向量};
 \node[conn] (t) at (5,1.15) {\ico{\faLanguage}\enspace\textbf{表达数据文本化}\\[.06cm]\scriptsize Cell2Sentence · C2S-Scale\\\scriptsize\color{Muted}排序基因名 $\rightarrow$ ``cell sentence''};
 \node[conn] (w) at (5,-1.15) {\ico{\faRobot}\enspace\textbf{工具调用}\\[.06cm]\scriptsize CellAgent · Biomni\\\scriptsize\color{Muted}研究任务 $\rightarrow$ 分析过程；本次不作主线};
 \draw[flow,line width=1.1pt] (p.east) -- node[cl,pos=.45]{作为骨干} ([yshift=.4cm]hub.west);
 \draw[flow] (m.east) -- node[cl,pos=.45]{提供先验} ([yshift=-.4cm]hub.west);
 \draw[flow,{Latex[length=2mm]}-{Latex[length=2mm]}] (t.west) -- node[cl,pos=.55]{文本接口} ([yshift=.4cm]hub.east);
 \draw[flow] (w.west) -- node[cl,pos=.55]{调用与分析} ([yshift=-.4cm]hub.east);
\end{tikzpicture}
\remark{\ico{\faStar}\enspace 本次关注预训练范式这一层：它直接产出可比较的扰动预测器，但每一个都需要与简单基线比较。}
\source{[5--7,14--21]。四层划分为本次调研的组织方式；各模型的编码与训练目标见附录。}
\end{frame}
